\begingroup\renewcommand{\arraystretch}{1.3}
\begin{longtable}{|>{\raggedright\arraybackslash}p{2.1cm}|>{\centering\arraybackslash}p{1.4cm}|>{\raggedright\arraybackslash}p{10.9cm}|}
\hline
\hdr{Date} & \hdr{Version} & \hdr{Description} \\ \hline
\endfirsthead
\hline
\hdr{Date} & \hdr{Version} & \hdr{Description} \\ \hline
\endhead
2026-09-29 & 1 & \textit{Innoweegbeg AI Reporting Tool --- Requirements Engineering} (Word): context and scope, four personas, three scenarios, use cases UC1--UC6, user stories US1--US10, FR1--FR16, NFR1--NFR12, traceability summary, workflow, test cases TC1--TC14 and the domain model of the ingestion code. \\ \hline
\rowcolor{light}
2026-10-09 & 1.1 & Review notes added: a consistency check against the research questions and the stakeholder survey. No corrections; three optional suggestions left to the author (kept unchanged in Section~\ref{sec:review}). \\ \hline
2026-10-30 & 2 & This PDF. The content of version 1.1 is kept word for word. Added: an implementation status for every use case, user story, requirement and test case, with the code or test that shows it; the domain model as now implemented; use cases UC7--UC9, user stories US11--US14, requirements FR17--FR21 and test cases TC15--TC22 for what was built since 29 September, each marked \added; a use-case diagram; the workflow redrawn as built; review notes of 30 October. \\ \hline
\end{longtable}
\endgroup
